% GENERATED FILE. Edit canonical lessons, then run npm run generate:books.
% canonical-source-hash: fnv1a64:f333da1b61a9e063
% canonical-lessons: FR-C527-fer-a-repasser, FR-C527-froisse, FR-C527-tremper, FR-C527-penderie, FR-C527-accrocher

\chapter{Repasser et ranger}
\label{ch:fr-repasser-et-ranger}
\hlchaptermodality{527}

\begin{chapteropening}
\textbf{By the end of this chapter:} \emph{I can talk about ironing and putting clothes away.}

Complete the last lesson of chapter 527: I can talk about ironing and putting clothes away.
\end{chapteropening}

\section[le fer à repasser]{le fer à repasser — iron (for clothes)}

\label{lesson:FR-C527-fer-a-repasser}



\begin{quote}
Before the new one: say the French for clothes peg, then the French for coat hanger.
\end{quote}

\subsection*{You'll want to know: le fer à repasser}
\textbf{le fer à repasser} — \textquotedblleft{}iron (for clothes)\textquotedblright{}.

\begin{quote}
Le fer à repasser est encore chaud.
\end{quote}

\begin{quote}
\emph{The iron is still hot.}
\end{quote}

Often \emph{le fer} alone. \emph{La planche à repasser} — the ironing board.

\begin{cousinweb}
\emph{le fer à repasser} is \textquotedblleft{}an iron for pressing\textquotedblright{}: \emph{fer} is inherited from Latin \emph{ferrum} \textquotedblleft{}iron\textquotedblright{} (Lewis \& Short), and \emph{repasser} is built on \emph{passer}.
\end{cousinweb}

\begin{grammarlens}[title={What you've built}]
One more word for telling, explaining and describing at B1.
\end{grammarlens}

\subsection*{Your turn}
\begin{itemize}
\raggedright
  \item \emph{Say it:} \emph{le fer à repasser}
  \item \emph{Say it:} \emph{le fer à repasser}, once more
  \item \emph{Say it:} a sentence of your own with \emph{le fer à repasser}
  \item \emph{Recall:} say the French for clothes peg, then the French for coat hanger, then say \emph{le fer à repasser} again
\end{itemize}

\begin{tcolorbox}[breakable,colback=teal!4,colframe=teal!35!black,title={Before you move on}]
What does le fer à repasser mean? (\textquotedblleft{}Iron (for clothes)\textquotedblright{}.) Say it once more.
\end{tcolorbox}
\section[froissé]{froissé — creased, crumpled; offended}

\label{lesson:FR-C527-froisse}



\begin{quote}
Before the new one: say the French for coat hanger, then the French for iron.
\end{quote}

\subsection*{You'll want to know: froissé}
\textbf{froissé} — \textquotedblleft{}creased, crumpled; offended\textquotedblright{}.

\begin{quote}
Ta chemise est toute froissée.
\end{quote}

\begin{quote}
\emph{Your shirt is all creased.}
\end{quote}

Also hurt, offended: \emph{il est froissé par ta remarque}.

\begin{cousinweb}
\emph{froissé} is the participle of \emph{froisser} \textquotedblleft{}to crumple\textquotedblright{} (EtymDB), from Vulgar Latin \emph{frustiāre} \textquotedblleft{}to break into bits\textquotedblright{}, from Latin \emph{frustum} \textquotedblleft{}a crumb, a bit\textquotedblright{} (Lewis \& Short).
\end{cousinweb}

\begin{grammarlens}[title={What you've built}]
One more word for telling, explaining and describing at B1.
\end{grammarlens}

\subsection*{Your turn}
\begin{itemize}
\raggedright
  \item \emph{Say it:} \emph{froissé}
  \item \emph{Say it:} \emph{froissé}, once more
  \item \emph{Say it:} a sentence of your own with \emph{froissé}
  \item \emph{Recall:} say the French for coat hanger, then the French for iron, then say \emph{froissé} again
\end{itemize}

\begin{tcolorbox}[breakable,colback=teal!4,colframe=teal!35!black,title={Before you move on}]
What does froissé mean? (\textquotedblleft{}Creased, crumpled; offended\textquotedblright{}.) Say it once more.
\end{tcolorbox}
\section[tremper]{tremper — to soak; to dip}

\label{lesson:FR-C527-tremper}



\begin{quote}
Before the new one: say the French for iron, then the French for creased.
\end{quote}

\subsection*{You'll want to know: tremper}
\textbf{tremper} — \textquotedblleft{}to soak; to dip\textquotedblright{}.

\begin{quote}
Je fais tremper le pantalon taché dans l'eau froide.
\end{quote}

\begin{quote}
\emph{I am soaking the stained trousers in cold water.}
\end{quote}

\emph{Faire tremper} — to leave to soak. \emph{Être trempé} — to be soaked.

\begin{cousinweb}
\emph{tremper} is inherited from Latin \emph{temperāre} \textquotedblleft{}to mix in due proportion\textquotedblright{} (Lewis \& Short), with its letters later shifted; the sense moved to mixing with water.
\end{cousinweb}

\begin{grammarlens}[title={What you've built}]
One more word for telling, explaining and describing at B1.
\end{grammarlens}

\subsection*{Your turn}
\begin{itemize}
\raggedright
  \item \emph{Say it:} \emph{tremper}
  \item \emph{Say it:} \emph{tremper}, once more
  \item \emph{Say it:} a sentence of your own with \emph{tremper}
  \item \emph{Recall:} say the French for iron, then the French for creased, then say \emph{tremper} again
\end{itemize}

\begin{tcolorbox}[breakable,colback=teal!4,colframe=teal!35!black,title={Before you move on}]
What does tremper mean? (\textquotedblleft{}To soak; to dip\textquotedblright{}.) Say it once more.
\end{tcolorbox}
\section[la penderie]{la penderie — wardrobe (for hanging clothes)}

\label{lesson:FR-C527-penderie}



\begin{quote}
Before the new one: say the French for creased, then the French for to soak.
\end{quote}

\subsection*{You'll want to know: la penderie}
\textbf{la penderie} — \textquotedblleft{}wardrobe (for hanging clothes)\textquotedblright{}.

\begin{quote}
Les manteaux sont dans la penderie de l'entrée.
\end{quote}

\begin{quote}
\emph{The coats are in the hall wardrobe.}
\end{quote}

A cupboard with a rail for hanging; \emph{l'armoire} has shelves too.

\begin{cousinweb}
\emph{la penderie} is built on \emph{pendre} \textquotedblleft{}to hang\textquotedblright{} with \emph{-erie}; \emph{pendre} is inherited from Latin \emph{pendēre} \textquotedblleft{}to hang\textquotedblright{} (Lewis \& Short).
\end{cousinweb}

\begin{grammarlens}[title={What you've built}]
One more word for telling, explaining and describing at B1.
\end{grammarlens}

\subsection*{Your turn}
\begin{itemize}
\raggedright
  \item \emph{Say it:} \emph{la penderie}
  \item \emph{Say it:} \emph{la penderie}, once more
  \item \emph{Say it:} a sentence of your own with \emph{la penderie}
  \item \emph{Recall:} say the French for creased, then the French for to soak, then say \emph{la penderie} again
\end{itemize}

\begin{tcolorbox}[breakable,colback=teal!4,colframe=teal!35!black,title={Before you move on}]
What does la penderie mean? (\textquotedblleft{}Wardrobe (for hanging clothes)\textquotedblright{}.) Say it once more.
\end{tcolorbox}
\section[accrocher]{accrocher — to hang (up), to hook}

\label{lesson:FR-C527-accrocher}



\begin{quote}
Before the new one: say the French for to soak, then the French for wardrobe.
\end{quote}

\subsection*{You'll want to know: accrocher}
\textbf{accrocher} — \textquotedblleft{}to hang (up), to hook\textquotedblright{}.

\begin{quote}
Il accroche sa veste derrière la porte.
\end{quote}

\begin{quote}
\emph{He hangs his jacket behind the door.}
\end{quote}

\emph{Accrocher un tableau} — to hang a picture. \emph{S'accrocher} — to hold on.

\begin{cousinweb}
\emph{accrocher} is built on \emph{croc} \textquotedblleft{}a hook\textquotedblright{} (EtymDB), which comes from Frankish \emph{krōk}, a Germanic word.
\end{cousinweb}

\begin{grammarlens}[title={What you've built}]
One more word for telling, explaining and describing at B1.
\end{grammarlens}

\subsection*{Your turn}
\begin{itemize}
\raggedright
  \item \emph{Say it:} \emph{accrocher}
  \item \emph{Say it:} \emph{accrocher}, once more
  \item \emph{Say it:} a sentence of your own with \emph{accrocher}
  \item \emph{Recall:} say the French for to soak, then the French for wardrobe, then say \emph{accrocher} again
\end{itemize}

\begin{tcolorbox}[breakable,colback=teal!4,colframe=teal!35!black,title={Before you move on}]
What does accrocher mean? (\textquotedblleft{}To hang (up), to hook\textquotedblright{}.) Say it once more.
\end{tcolorbox}
